% Generated draft: requires mathematical and source review.
\documentclass[11pt]{article}
\usepackage[margin=25mm]{geometry}
\usepackage{fontspec}
\setmainfont{Latin Modern Roman}
\usepackage{amsmath,amssymb,mathtools,unicode-math}
\setmathfont{Latin Modern Math}
\usepackage{enumitem,needspace,xurl,hyperref}
\setlength{\parindent}{0pt}
\setlength{\parskip}{5pt}
\setlength{\emergencystretch}{3em}
\newcommand{\N}{\mathbb{N}}
\newcommand{\Z}{\mathbb{Z}}
\newcommand{\Q}{\mathbb{Q}}
\newcommand{\R}{\mathbb{R}}
\newcommand{\C}{\mathbb{C}}
\newcommand{\F}{\mathbb{F}}
\newcommand{\A}{\mathbb{A}}
\newcommand{\PP}{\mathbb{P}}
\newcommand{\E}{\mathbb{E}}
\newcommand{\eps}{\varepsilon}
\newcommand{\dd}{\,\mathrm{d}}
\begin{document}
\section*{UnsolvedMath ID: 20000636}
\textbf{Source problem code:} catalogue entry
\subsection{Definitions and mathematical statement}
Yan Fyodorov's source asks whether the distribution of real zeros of a random entire function, real on $\R$ and of order one, determines the distribution of all zeros of $f$. The extracted wording is ``on order 1''; interpreting this as ``of order 1'' follows the supplied verification note. That note also says there is no evidence for a prime on the final $f$, so this draft concerns $f$, not $f'$.

For a nonzero entire function let $Z_f=\sum_{f(\zeta)=0}\operatorname{mult}_\zeta(f)\,\delta_\zeta$ be its zero divisor, and let $Z_f^{\R}$ be its restriction to $\R$. With $M_f(r)=\max_{|z|=r}|f(z)|$, order one means $\limsup_{r\to\infty}\log\log M_f(r)/\log r=1$. For a function of finite exponential type, its type is $\limsup_{r\to\infty}r^{-1}\log M_f(r)$ and its indicator is $h_f(\theta)=\limsup_{r\to\infty}r^{-1}\log|f(re^{i\theta})|$. The Cartwright class consists of entire functions of finite exponential type with
\[
 \int_{\R}\frac{\log^+|f(x)|}{1+x^2}\,dx<\infty.
\]
The original statement specifies no law for the random entire function beyond the given real-zero distribution.

The supplied aid records the following nonuniqueness family. Let $U,V$ be independent uniform variables on $[0,1)$, and, for $a>0$, set
\[
 F_a(z)=\sin\bigl(\pi(z-V)\bigr)
        \sin\bigl(\pi(z-U-ia)\bigr)
        \sin\bigl(\pi(z-U+ia)\bigr).
\]
It reports that each $F_a$ is real entire, belongs to the Cartwright class, has exact order one and type $3\pi$, and has indicator
\[
 h_{F_a}(\theta)=3\pi|\sin\theta|.
\]
Its real zeros are $V+\Z$, while its nonreal zeros are $U+\Z+ia$ and $U+\Z-ia$, with their usual multiplicities. Thus every $a$ gives the same real-zero-divisor law, but distinct positive values of $a$ give mutually singular nonreal-zero laws. Both the real and complete zero-divisor processes are stationary under real translations; this assertion concerns zero processes and does not require stationarity of the random functions themselves.

The aid also reports that finite conjugate-pair polynomial insertion can preserve a prescribed finite Taylor jet as well as the real divisor, order, type, and indicator. Precisely, for a fixed jet order $k$ at the origin, multiplication by a real polynomial $P$ having no real zero and satisfying $P(z)=1+O(z^{k+1})$ preserves $f^{(j)}(0)$ for $0\leq j\leq k$ and leaves $Z_f^{\R}$ unchanged, while inserting the nonreal zeros of $P$. The preservation of the order-one growth data is the reported finite-polynomial effect; stationarity is not asserted for this insertion construction. These examples show that a unique full zero law cannot be inferred from the stated real-zero data, even with the stronger fixed growth restrictions in the stationary family. They do not classify which full zero laws are possible under additional ensemble assumptions.
\subsection{Short English statement}
A real-zero distribution alone does not specify all complex zeros. The supplied family has the same stationary real lattice of zeros and the same order, type and indicator for every parameter, but its nonreal zeros lie at parameter-dependent heights. Even a prescribed finite Taylor jet can remain unchanged when finitely many conjugate nonreal zeros are inserted.
\subsection{Sources}
\begin{itemize}
\item Local UnsolvedMath registry, permanent ID $20000636$, source code \texttt{catalogue entry}, supplied complete source snapshot, verification note and research summary.
\item Yan Fyodorov, Question 11, AIM workshop \emph{Random analytic functions}, supplied source URL \url{https://aimath.org/WWN/randomzeros/arcc1.pdf}; canonical local record \texttt{aim-analysis-notes.json}, \texttt{notes[167]}. This citation does not represent a new PDF visit.
\item The original literature field is empty; no authoritative publication reference for the counterexample family or insertion statement is supplied.
\item Research attempt 1: partial result; problem partially solved at the recorded run; rigor label ``proved'', novelty confidence ``low''. The machine-generated aid is not peer reviewed and requires independent verification before citation.
\end{itemize}
\end{document}
